\section{Related Work}
\subsection{Atmospheric Forecast Post-processing}
Station-specific model output statistics is established. Riccio and Chianese
study statistical improvement of CAMS with spatial verification
\cite{riccio2024mos}; Petetin et al. examine the trade-off between continuous
and categorical ozone-forecast skill \cite{petetin2022mos}. Operational
regional CAMS production also includes statistical post-processing
\cite{colette2025cams}. Consequently, introducing a correction learner alone
does not establish novelty. AirGuard's contribution under investigation is
the audited Indian station evaluation and its warning behaviour.

\subsection{Uncertainty and Warning Decisions}
Conformal methods provide a framework for data-driven prediction intervals.
Time-series and distribution-shift settings require attention to dependence
and changing coverage \cite{xu2021enbpi,gibbs2021aci}. A narrow interval or
near-nominal aggregate coverage is insufficient if high-pollution categories
are systematically missed. Proper forecast evaluation also separates the
quantity predicted from the decision made from it \cite{gneiting2007scoring}.
Here, vulnerability groups specify exposure-warning tolerances; they do not
constitute clinical cohorts or individual health-risk models.

\subsection{Scope of the Comparison}
Published methods use different targets, locations and verification periods.
Their published scores are therefore not inserted into AirGuard's numerical
leaderboard. The shared-observation CAMS comparison, spatial inference
regimes and uncertainty failures must be established with this study's own
artifacts. A complete novelty assessment remains pending the primary study.
